% !TEX program = lualatex
% ECON 201, Worksheet 20: Taxes and Price Ceilings.
% Compile with LuaLaTeX so the PDF is tagged. Build from the course root:
%   latexmk -lualatex -cd worksheets/worksheet20.tex && latexmk -c -cd worksheets/worksheet20.tex
% Two sides, one sheet. Every number here is computed by
% slides/figures/make_lecture20.py. Front: gasoline at stations in town, demand
% P = 8 - Q/10, supply P = 2 + Q/10, taxed supply P = 3 + 3Q/20. Back: apartments
% near campus, demand P = 2100 - 3Q, supply P = 300 + 3Q, a rent cap of $900.
% The script also draws figures/gasoline.pdf and figures/apartments.pdf.
% No definitions panel: nothing here needs one.
\DocumentMetadata{
  pdfstandard=UA-2,
  pdfversion=2.0,
  lang=en-US,
  tagging=on
}
\documentclass[11pt]{article}

\usepackage[letterpaper, margin=0.75in, top=0.5in, bottom=0.45in]{geometry}
\usepackage{fontspec}
\setmainfont{Fira Sans}
\usepackage{amsmath}
\usepackage{amssymb}
\usepackage{graphicx}
\usepackage{booktabs}
\usepackage{array}
\usepackage{xcolor}
\usepackage{needspace}
\usepackage[most]{tcolorbox}
\definecolor{cream}{HTML}{FDF3EA}
\definecolor{accent}{HTML}{BF5700}
\usepackage{titlesec}
\titleformat{\section}{\large\bfseries\color{accent}}{}{0pt}{}
\titlespacing*{\section}{0pt}{4pt}{2pt}
\usepackage{hyperref}
\hypersetup{pdftitle={ECON 201 Worksheet: Taxes and Price Ceilings}, pdfauthor={Div Bhagia}, hidelinks}
\pagestyle{empty}
\setlength{\parindent}{0pt}
\setlength{\parskip}{2pt}

% Ruled answer space: n lines. Plain TeX loop; pgffor's \foreach breaks the
% enumerate counter when used inside an item.
\newcount\answerlinecount
\newcommand{\answerlines}[1]{%
  \par\vspace{1pt}\answerlinecount=#1
  \loop\ifnum\answerlinecount>0
    \rule{\linewidth}{0.3pt}\par\vspace{5pt}%
    \advance\answerlinecount by -1
  \repeat
  \vspace{-5pt}}

\begin{document}

{\LARGE\bfseries\color{accent} Worksheet: Taxes and Price Ceilings}\hfill{\large ECON 201}

\vspace{4pt}
\section{Gasoline at stations in town}

Many stations in town sell the same regular gasoline. In the graph, the solid
lines show the original demand and supply curves.

\begin{center}
\includegraphics[width=0.74\textwidth, alt={A plotting grid, price in dollars a
gallon on the vertical axis from 0 to 8 with gridlines every dollar, and
thousand gallons of gasoline a day on the horizontal axis from 0 to 60 with
gridlines every 10. The demand curve falls in a straight line from 8 dollars at
zero gallons to 2 dollars at 60 thousand. The supply curve rises in a straight
line from 2 dollars at zero gallons to 8 dollars at 60 thousand. They cross at
30 thousand gallons and 5 dollars. A dashed line labelled Supply with tax runs
50 percent above supply, rising from 3 dollars at zero gallons to 8 dollars, and
crosses demand at 20 thousand gallons and 6 dollars.}]{figures/gasoline.pdf}
\end{center}

\vspace{4pt}
1. What are the equilibrium price and quantity? Mark this point, A, on the graph.
\answerlines{1}

2. The city puts a 50\% tax on sellers, so their marginal cost of supplying each
gallon is 50\% higher: the dashed line shows supply with the tax. How many gallons
are sold, what do buyers pay, and what do sellers keep?
\answerlines{2}

3. What share of the tax do buyers bear, and what share do sellers bear?
\answerlines{2}

\clearpage
\section{Apartments near campus}

Many landlords near campus rent out similar apartments. In the graph, the lines
show the demand and supply curves.

\begin{center}
\includegraphics[width=0.74\textwidth, alt={A plotting grid, rent in dollars a
month on the vertical axis from 0 to 2,100 with gridlines every 300 dollars, and
the number of apartments on the horizontal axis from 0 to 600 with gridlines
every 100. The demand curve falls in a straight line from 2,100 dollars at zero
apartments to 300 dollars at 600. The supply curve rises in a straight line from
300 dollars at zero apartments to 2,100 dollars at 600. They cross at 300
apartments and 1,200 dollars.}]{figures/apartments.pdf}
\end{center}

\vspace{4pt}
4. What are the equilibrium rent and quantity? Mark this point, A, on the graph.
\answerlines{1}

5. The city caps the rent at \$900 a month. How many apartments do landlords
offer, how many do renters want, and how many are rented?
\answerlines{2}

6. Do the apartments go to the renters who value them most? Explain.
\answerlines{2}

\vspace{10pt}
\emph{Practice at home}: for the gasoline tax, calculate the consumer surplus,
producer surplus, tax revenue, and deadweight loss. Do the same for the rent cap:
consumer surplus, producer surplus, and deadweight loss.

\end{document}
